% =====================================================================================
% Insert addressing the estimand objection (REVIEW.md R3): Eq. (9) is written over inputs,
% the measurement is indexed by behaviour.
%
% ALL NUMBERS ARE MEASURED. No placeholders.
%   Source : results/hpc_vicuna_autodan/gold.jsonl -> fusion/cross_breach.csv
%   Script : scripts/input_level_phi.py, input built by scripts/10_build_cross_breach.py
%   Data   : fusion/r3_estimand_comparison.csv, fusion/run_r3.log
%
% Place immediately after the \paragraph{Setup.} block of \S\ref{sec:corr-measure}, or as
% the first subsection after it. The closing paragraph is written for \S\ref{sec:limits-corr}.
% =====================================================================================

\subsection{Which estimand does \texorpdfstring{$\phi$}{phi} measure?}
\label{sec:estimand}

Eq.~\eqref{eq:joint} is written over inputs: $b_l(i)$ is the event that layer $l$ is breached
on input $i$, and the joint residual is derived with both layers facing the same input. Our
measurement is indexed by \emph{behaviour} rather than by input, because the adaptive attack
is re-optimised against each defense separately. For a pair $(d_1,d_2)$, $d_1$'s entry comes
from a prompt optimised against $d_1$ and $d_2$'s from a different prompt optimised against
$d_2$. These are two distinct quantities and the distinction should be made explicitly rather
than left to the reader:

\begin{itemize}
\item $\phi_{\mathrm{beh}}$ scores each defense on prompts optimised against itself. This is
what a deployer faces when the adversary is free to retarget after observing which layer
blocked, and it is the quantity reported in Table~\ref{tab:dependency}.
\item $\phi_{\mathrm{inp}}$ scores both defenses on the \emph{same} prompt. This is the
quantity Eq.~\eqref{eq:joint} is written over.
\end{itemize}

They coincide only if the attack transfers perfectly between defenses, which is precisely
what \S\ref{par:attack-class} shows to be false. We therefore measure the difference rather
than assume it away.

\paragraph{A shared-input estimate from data already collected.}
The two static baselines are shared-input by construction: the plain behaviour prompt is
byte-identical across defenses, and the transfer suffix is a single fixed string appended to
every prompt, so neither consults the defense. Computing $\phi$ within each of those prompt
sets gives $\phi_{\mathrm{inp}}$ directly, with no new attack search and no new generations.

\begin{center}
\begin{tabular}{lccc}
\toprule
Pair & $\phi_{\mathrm{beh}}$ & $\phi_{\mathrm{inp}}$ (plain) & $\phi_{\mathrm{inp}}$ (transfer) \\
\midrule
probe$_{16}$ $\times$ probe$_{8}$   & 0.75 & 0.86 & 0.90 \\
perplexity $\times$ probe$_{8}$     & 0.62 & 0.92 & --   \\
perplexity $\times$ probe$_{16}$    & 0.59 & 0.93 & --   \\
perplexity $\times$ refusal-prime   & 0.58 & 0.64 & --   \\
refusal-prime $\times$ SmoothLLM    & 0.56 & 0.20 & 0.58 \\
probe$_{8}$ $\times$ SmoothLLM      & 0.48 & 0.29 & 0.33 \\
probe$_{16}$ $\times$ token-anomaly & 0.32 & 0.93 & 0.85 \\
perplexity $\times$ token-anomaly   & 0.35 & 1.00 & --   \\
\bottomrule
\end{tabular}
\end{center}

\noindent
Across all fifteen pairs measurable on the plain set, $\phi_{\mathrm{inp}}$ has the
\emph{same sign} as $\phi_{\mathrm{beh}}$ in $15/15$ cases and is positive in $15/15$, with a
mean shift of $+0.14$. On the transfer set, where ten pairs are measurable, the agreement is
$10/10$ on sign, $10/10$ positive, mean shift $+0.10$, and every one of the ten bootstrap
intervals excludes zero. The conclusion of \S\ref{sec:corr-measure} therefore holds under
both estimands, and the behaviour-level figure we report is if anything the
\emph{conservative} one: measured on a common input, these defenses fail together more
tightly, not less.

That the input-level correlation is the larger of the two is what one should expect. Holding
the prompt fixed removes the only source of variation that could decorrelate the two layers,
namely the adversary's freedom to construct a different input for each. Allowing that freedom
is what makes $\phi_{\mathrm{beh}}$ the operationally relevant quantity for a deployer, and
it is why we report it as the headline.

\paragraph{What this estimate does not settle.}
The static baselines are weak attacks, with marginals between $0.03$ and $0.28$ against
between $0.35$ and $0.68$ for the fluent adversary. Several plain-set pairs therefore rest on
six to eight breach events: five of the fifteen have bootstrap intervals spanning zero, all
of them involving SmoothLLM, and the perplexity $\times$ token-anomaly value of $1.00$ means
only that the same seven behaviours breached both filters. These are brackets, not
replacements. The definitive version is a cross-evaluation matrix in which the prompts
optimised against each defense are scored against all the others; that is a scoring pass over
prompts that already exist rather than a new attack search, and we identify it as the single
most valuable extension of this experiment alongside the larger behaviour set of
\S\ref{sec:confound}.

% =====================================================================================
% APPEND TO \S\ref{sec:limits-corr} (Limitations)
% =====================================================================================

\textbf{Estimand.} The reported $\phi$ is behaviour-level: each defense is scored on prompts
optimised against itself, which matches a retargeting adversary but not the same-input
quantity Eq.~\eqref{eq:joint} is written over. \S\ref{sec:estimand} bounds the difference
using the two shared-input baselines and finds agreement in sign on all measurable pairs,
with the input-level estimate the larger of the two, but those baselines are weak attacks and
several of their pairs are underpowered. A full cross-evaluation under the fluent adversary
remains outstanding.
